% ==============================================================================
% CAPÍTULO 1 - INTRODUÇÃO
% ==============================================================================

\chapter{Introdução}
\label{cap:introducao}

% Contextualização do tema

Nesta seção, você deve descrever o cenário de tecnologia e a empresa onde o estágio foi realizado. 
Comece contextualizando a grande área do conhecimento (como a Computação ou Engenharia de Software)
 e relacione-a com a atuação da empresa. Depois, fale sobre a área mais específica do seu trabalho 
 (ex: desenvolvimento web, banco de dados ou infraestrutura). Por fim, descreva o problema ou necessidade
  que existia no contexto da empresa e a solução que você propôs ou ajudou a desenvolver. O texto deve começar 
falando da area e mercado em geral e terminar focando no seu papel específico, apresentando o 
objetivo deste relatório: registrar as atividades que você fez e o que você aprendeu.



%-------------------------------------------------
% Justificativa
%-------------------------------------------------
\section{Justificativa}
\label{sec:justificativa}

A justificativa deve evidenciar a importância da experiência prática para a 
consolidação das competências previstas no currículo acadêmico de Computação/Engenharia de Software. 
É necessário destacar como a resolução de problemas reais contribui para o amadurecimento profissional 
e técnico do estagiário. Além disso, mencione o valor agregado para a empresa 
(ex: otimização de processos, desenvolvimento de módulos, suporte técnico) e 
como essa integração universidade-empresa fortalece a formação acadêmica.

%-------------------------------------------------
% Objetivos
%-------------------------------------------------
\section{Objetivos}
\label{sec:objetivos}

\subsection{Objetivo Geral}
\label{subsec:objective-geral}

O objetivo geral deve expressar a meta principal do estágio de forma sintetizada.
Ele deve iniciar com um verbo no infinitivo que denote ação e resultado. Exemplos comuns na nossa área:\\
 - Desenvolver sistemas de software utilizando tecnologias X e Y...\\
 - Implementar rotinas de automação para...\\
 - Analisar e otimizar a infraestrutura de rede da empresa...

\subsection{Objetivos Específicos}
\label{subsec:objetivos-especificos}

Os objetivos específicos devem detalhar as etapas técnicas e operacionais que permitirão o alcance do objetivo geral. 
Devem ser mensuráveis e diretos. Exemplos comuns na nossa área:

\begin{alineas}
    \item Realizar o levantamento de requisitos junto aos usuários ou gestores;
    \item Elaborar a documentação técnica e diagramas de arquitetura (UML, MER);
    \item Codificar módulos ou funcionalidades utilizando a linguagem [Java/TypeScript/etc.];
    \item Executar testes unitários e de integração para garantir a qualidade do software;
    \item Documentar os manuais de uso ou processos de implantação.
\end{alineas}




%-------------------------------------------------
% Estrutura do Trabalho
%-------------------------------------------------
\section{Estrutura do Trabalho}
\label{sec:estrutura}
Nesta seção, você deve explicar como o seu relatório está organizado.
Cite cada capítulo do trabalho (Introdução, Empresa, Metodologia, Revisão, Desenvolvimento e Conclusão) e 
resuma em uma frase curta o que o leitor encontrará em cada um deles. 
O objetivo é servir como um "guia" que antecipa a ordem das informações que serão apresentadas.

Exemplo:

Para melhor compreensão, este relatório está estruturado em capítulos. Além desta Introdução,
o \autoref{ch:apresentacao-empresa} apresenta o perfil da organização onde o estágio foi realizado. 
O \autoref{cap:metodologia} detalha os procedimentos, métodos ....